\chapter{Completitud de los números reales}

\section{Introducción: La Propiedad de Completitud}

Hemos llegado ahora a una de las propiedades más fundamentales de los números
reales: la \emph{completitud}. Esta es la propiedad que distingue a los números
reales de los números racionales y hace posible el análisis.

Sabemos que las sucesiones de Cauchy de números racionales convergen a números
reales. De hecho, así es como construimos los números reales: como clases de
equivalencia de sucesiones de Cauchy de racionales. Registremos esto como un
teorema:

\begin{theorem}[Conv\_of\_IsCauchy]
Si una sucesión $a : \mathbb{N} \to \mathbb{Q}$ satisface \texttt{IsCauchy a},
entonces \texttt{SeqConv a} se cumple, es decir, la sucesión converge.
\end{theorem}

Dado que dicho límite existe, podemos definirlo explícitamente:

\begin{definition}[Real\_of\_CauSeq]
Toma \texttt{IsCauchy a} (para $a : \mathbb{N} \to \mathbb{Q}$) y devuelve el
número real al que converge $a$.
\end{definition}

Este número real hace lo que queremos:

\begin{theorem}[SeqLim\_of\_Real\_of\_Cau]
Si $\mathit{ha} : \texttt{IsCauchy}\ (a : \mathbb{N} \to \mathbb{Q})$, entonces
$\texttt{SeqLim}\ a\ (\texttt{Real\_of\_CauSeq}\ \mathit{ha})$ se cumple, es
decir, la sucesión $a$ converge al número real definido por
\texttt{Real\_of\_CauSeq}.
\end{theorem}

Pero ahora debemos preguntarnos: ¿qué ocurre con las sucesiones de Cauchy de
números reales? ¿Necesitamos aún otro sistema numérico (los «hiperreales» o quizás
los «sureales», para los aficionados a Conway\ldots) para capturar los límites de
los reales?

\noindent ¡No! Y esto es un fenómeno general: cuando se «completa» un espacio
reemplazándolo por clases de equivalencia de sucesiones de Cauchy de ese espacio,
el espacio mismo queda automáticamente «completo». Es decir, todas las sucesiones
de Cauchy en el espacio completado ya convergen dentro del mismo espacio.
¡Demostrémoslo para los reales!

\section{La Pregunta: ¿Son Completos los Reales?}

Queremos demostrar que toda sucesión de Cauchy de números reales converge a un
número real.

Pero acabamos de aprender que los números reales son, en sí mismos, sucesiones de
Cauchy de racionales (o más exactamente, clases de equivalencia de las mismas).
Ahora que los números reales están representados por sucesiones de Cauchy de
racionales, ¿qué significa que $(x_n)$ sea una sucesión de Cauchy de reales?
¿¡¿Una sucesión de Cauchy de sucesiones de Cauchy?!?

Desempaquemos eso.

\subsection{El Cuadro del Arreglo Doble}

Digamos que tenemos una sucesión infinita de números reales,
$x : \mathbb{N} \to \mathbb{R}$, que pensamos como $x = (x_0, x_1, x_2, \ldots)$.

Cada real $x_n$ está representado en secreto por una sucesión de Cauchy
$(q_{n,0},\, q_{n,1},\, q_{n,2},\, \ldots)$ de racionales. Así que cada $x_n$ es
una función $\mathbb{N} \to \mathbb{Q}$, y la $x$ original es en realidad una
función:
\[
  q : \mathbb{N} \to \mathbb{N} \to \mathbb{Q}
\]
Es decir, $x$ es realmente un arreglo de racionales con doble índice:

\[
\begin{array}{c|ccccl}
  & j=0 & j=1 & j=2 & j=3 & \cdots \\
\hline
i=0: & q_{0,0} & q_{0,1} & q_{0,2} & q_{0,3} & \cdots \quad \leftarrow \text{representa } x_0 \\
i=1: & q_{1,0} & q_{1,1} & q_{1,2} & q_{1,3} & \cdots \quad \leftarrow \text{representa } x_1 \\
i=2: & q_{2,0} & q_{2,1} & q_{2,2} & q_{2,3} & \cdots \quad \leftarrow \text{representa } x_2 \\
i=3: & q_{3,0} & q_{3,1} & q_{3,2} & q_{3,3} & \cdots \quad \leftarrow \text{representa } x_3 \\
\vdots & \vdots & \vdots & \vdots & \vdots & \ddots
\end{array}
\]

Cada fila es una sucesión de Cauchy de racionales $q_{n,m}$ que converge al real
correspondiente $x_n$. ¡Asombroso! Llamemos $\mathit{hq}$ a la condición de que
cada una de estas filas sea una sucesión de Cauchy:
\[
  \mathit{hq} : \forall n,\; \texttt{IsCauchy}\ (q\; n)
\]

Entonces decir que esta sucesión (de sucesiones) $x_n$ es en sí misma «de Cauchy»
significa (como ya sabemos) que:
\[
  \forall \varepsilon > 0,\; \exists N,\; \forall n \geq N,\; \forall m \geq n,\;
  |\texttt{Real\_of\_CauSeq}\ (\mathit{hq}\; m) - \texttt{Real\_of\_CauSeq}\
  (\mathit{hq}\; n)| < \varepsilon
\]

\subsection{El Objetivo}

Ahora, supongamos esta condición sobre nuestra sucesión de Cauchy (de sucesiones)
$x$. Nuestro objetivo es demostrar que esta sucesión tiene un límite real; es decir,
que existe una única sucesión $y : \mathbb{N} \to \mathbb{Q}$ que satisface
\[
  \mathit{hy} : \texttt{IsCauchy}\ y
\]
y tal que $x_n \to y$, lo que significa:
\[
  \texttt{SeqLim}\ (\lambda n \mapsto \texttt{Real\_of\_CauSeq}\ (\mathit{hq}\; n))\
  (\texttt{Real\_of\_CauSeq}\ \mathit{hy})
\]
¿¡¿Cómo se supone que debemos construir tal cosa?!?

\section{La Idea Clave: La Diagonalización de Cantor}

La idea clave recuerda a los argumentos de diagonalización de Cantor, muy
presentes en el análisis real: ¡necesitamos extraer una única sucesión de racionales
de este arreglo!

\subsection{El Intento Ingenuo}

\textbf{El intento ingenuo:} Tomar la sucesión diagonal
$(q_{0,0},\, q_{1,1},\, q_{2,2},\, q_{3,3},\, \ldots)$.

\textbf{El problema:} ¡Cada fila converge a su propio ritmo! La fila $k$ podría
necesitar llegar hasta el índice $1000k$ antes de estar dentro de $\varepsilon$ de
su límite, y la fila $k+1$ podría necesitar llegar hasta el índice $k^2$, etc. Pero
la diagonal solo muestrea la fila $k$ en la posición $k$, ¡lo cual puede ser
demasiado pronto!

\subsection{La Solución: Elegir Puntos de Convergencia}

\textbf{La solución:} Para cada fila $k$, necesitamos escoger un índice $N(k)$
suficientemente alejado para que la fila $k$ haya «convergido lo suficiente» en ese
punto. Luego tomamos:
\[
  y(k) = q_{k,\, N(k)}
\]
¡Esta sucesión $y$ será de Cauchy y representará el límite de nuestra sucesión de
reales!

¿Cómo elegir $N(k)$? Usamos el hecho de que cada $x_k$ (como número real) es el
límite de su sucesión de Cauchy. Entonces, para cualquier tolerancia (digamos
$1/k$, o más exactamente $1/(k+1)$ para que siempre sea estrictamente positiva),
existe un índice $N(k)$ tal que para todo $m \geq N(k)$, $q_{k,m}$ permanece dentro
de $1/(k+1)$ de $x_k$ (es decir, de $q_{k,n}$ para $n \geq N(k)$). ¡Eso es lo que
elegiremos!

\section{La Estrategia de Demostración}

Nuestra tarea es:
\begin{enumerate}
  \item Construir la sucesión $y$ eligiendo índices apropiados $N(k)$ de cada fila.
  \item Demostrar que $y$ es de Cauchy.
  \item Demostrar que la sucesión de reales converge al real representado por $y$.
\end{enumerate}

Recorramos la demostración formal.

\subsection{Paso 1: Elegir los Índices}

Para cada $n : \mathbb{N}$, usamos \texttt{SeqLim\_of\_Real\_of\_Cau}
$(\mathit{hq}\; n)$ con tolerancia $1/(n+1)$ para obtener un índice $N(n)$ tal que:
\[
  \forall m \geq N(n),\quad |q_{n,m} - x_n| < \frac{1}{n+1}
\]
Definimos:
\[
  y(n) = q_{n,\, N(n)}
\]
Esto nos da una única sucesión de racionales, uno de cada fila, elegido en puntos
donde esa fila ya ha convergido suficientemente cerca de su límite.

\subsection{Paso 2: Demostrar que $y$ es de Cauchy}

Para demostrar que $y$ es de Cauchy, necesitamos demostrar que para todo
$\varepsilon > 0$, existe $N_0$ tal que para todo $n, m \geq N_0$, se tiene
$|y_m - y_n| < \varepsilon$.

La idea clave es un argumento de desigualdad triangular:
\[
  |y_m - y_n| = |q_{m,N(m)} - q_{n,N(n)}|
  \leq |q_{m,N(m)} - x_m| + |x_m - x_n| + |x_n - q_{n,N(n)}|
\]

\textbf{Acotando el primer y el tercer término:} Necesitamos que
$|q_{m,N(m)} - x_m|$ y $|q_{n,N(n)} - x_n|$ sean pequeños. Por nuestra elección de
$N(m)$ y $N(n)$, tenemos:
\[
  |q_{m,N(m)} - x_m| < \frac{1}{m+1}, \qquad
  |q_{n,N(n)} - x_n| < \frac{1}{n+1}
\]
Para que sean menores que $\varepsilon/3$, necesitamos $1/(m+1) < \varepsilon/3$ y
$1/(n+1) < \varepsilon/3$. Por la propiedad arquimediana, existe $N_2$ tal que
$1/(N_2+1) < \varepsilon/3$. Para todo $n, m \geq N_2$, estas cotas se cumplen.

\textbf{Acotando el término del medio:} Necesitamos que $|x_m - x_n| < \varepsilon/3$.
Como $x$ es una sucesión de Cauchy de reales, existe $N_1$ tal que para todo
$m, n \geq N_1$, se tiene $|x_m - x_n| < \varepsilon/3$.

\textbf{Combinando:} Sea $N_0 = N_1 + N_2$. Entonces para todo $m, n \geq N_0$, los
tres términos están acotados por $\varepsilon/3$, así que:
\[
  |y_m - y_n| < \frac{\varepsilon}{3} + \frac{\varepsilon}{3} + \frac{\varepsilon}{3} = \varepsilon
\]
Por lo tanto, $y$ es de Cauchy.

\subsection{Paso 3: Demostrar que $x_n \to \texttt{Real\_of\_CauSeq}\ \mathit{hy}$}

Sea $L = \texttt{Real\_of\_CauSeq}\ \mathit{hy}$ el número real representado por
$y$. Necesitamos demostrar que $x_n \to L$, es decir:
\[
  \forall \varepsilon > 0,\; \exists N_0,\; \forall n \geq N_0,\quad |x_n - L| < \varepsilon
\]
De nuevo, usamos una desigualdad triangular:
\[
  |x_n - L| = |x_n - y_n + y_n - L| \leq |x_n - y_n| + |y_n - L|
\]

\textbf{Acotando el primer término:} Tenemos $y_n = q_{n,N(n)}$, así que:
\[
  |x_n - y_n| = |x_n - q_{n,N(n)}| < \frac{1}{n+1}
\]
Para que sea menor que $\varepsilon/2$, necesitamos $1/(n+1) < \varepsilon/2$. Por
la propiedad arquimediana, existe $N_2$ tal que para todo $n \geq N_2$, esto se
cumple.

\textbf{Acotando el segundo término:} Como $y$ es de Cauchy con límite $L$, por
\texttt{SeqLim\_of\_Real\_of\_Cau} $\mathit{hy}$, existe $N_3$ tal que para todo
$n \geq N_3$:
\[
  |y_n - L| < \frac{\varepsilon}{2}
\]

\textbf{Un ingrediente más:} También necesitamos usar el hecho de que $x$ en sí
misma es de Cauchy. Existe $N_1$ tal que para todo $n, m \geq N_1$,
$|x_n - x_m| < \varepsilon/2$.

\textbf{Combinando:} Sea $N_0 = N_1 + N_2 + N_3$. Entonces para todo $n \geq N_0$:
\[
  |x_n - L| \leq |x_n - y_n| + |y_n - L| < \frac{\varepsilon}{2} + \frac{\varepsilon}{2} = \varepsilon
\]
Por lo tanto, $x_n \to L$, completando la demostración.

\section{La Demostración en Lean}

\begin{minted}[fontsize=\small, breaklines, linenos]{lean}
Statement Reals_are_Complete (q : N → N → Q) (x : N → R)
    (hq : ∀ n, IsCauchy (q n))
    (hx : ∀ n, x n = Real_of_CauSeq (hq n))
    (hxCau : IsCauchy x) :
    ∃ (y : N → Q) (hy : IsCauchy y),
      SeqLim x (Real_of_CauSeq hy) := by
  choose N hN using fun (n : N) ↦
    SeqLim_of_Real_of_Cau (hq n) (1 / (n + 1)) (by bound)
  have hN : ∀ n, ∀ m ≥ N n, |(q n m) - x n| < 1 / (n + 1) := by
    intro n m hm
    rewrite [hx n]
    apply hN n m hm
  let y := fun n ↦ q n (N n)
  use y
  have hy : IsCauchy y := by
    intro ε hε
    have hε' : (0 : R) < ε := by exact_mod_cast hε
    choose N1 hN1 using hxCau (ε / 3) (by norm_num; bound)
    choose N2 hN2 using ArchProp
      (show (0 : R) < ε / 3 by norm_num; bound)
    use N1 + N2
    intro n hn m hm
    change |(q m (N m)) - q n (N n)| < ε
    specialize hN1 n (by bound) m hm
    have qnBnd := hN n (N n) (by bound)
    have qmBnd := hN m (N m) (by bound)
    have f1 : |(q m (N m) : R) - q n (N n)| =
        |(q m (N m) - x m) + ((x n - q n (N n)) + (x m - x n))| :=
      by ring_nf
    have f2 : |(q m (N m) - x m) +
        ((x n - q n (N n)) + (x m - x n))| ≤
        |(q m (N m) - x m)| + |(x n - q n (N n)) + (x m - x n)| :=
      by apply abs_add
    have f3 : |(x n - q n (N n)) + (x m - x n)| ≤
        |(x n - q n (N n))| + |(x m - x n)| := by apply abs_add
    have f3' : |(x n - q n (N n))| = |q n (N n) - x n| :=
      by apply abs_sub_comm
    field_simp at hN2
    have hn' : (N1 : R) + N2 ≤ n := by exact_mod_cast hn
    have hm' : (n : R) ≤ m := by exact_mod_cast hm
    have f4' : (N2 : R) ≤ n := by bound
    have f4'' : (N2 : R) * ε ≤ n * ε := by bound
    have f4 : (1 : R) / (n + 1) < ε / 3 := by field_simp; bound
    have f5' : (N2 : R) ≤ m := by bound
    have f5'' : (N2 : R) * ε ≤ m * ε := by bound
    have f5 : (1 : R) / (m + 1) < ε / 3 := by field_simp; bound
    have f6 : |(q m (N m) : R) - q n (N n)| < ε := by
      linarith [f1, f2, f3, f4, f5, qnBnd, qmBnd, hN1, f3']
    exact_mod_cast f6
  use hy
  intro ε hε
  choose N1 hN1 using hxCau (ε / 2) (by norm_num; bound)
  choose N2 hN2 using ArchProp
    (show (0 : R) < ε / 2 by norm_num; bound)
  choose N3 hN3 using SeqLim_of_Real_of_Cau hy (ε / 2)
    (by norm_num; bound)
  use N1 + N2 + N3
  intro n hn
  let L := Real_of_CauSeq hy
  change |x n - L| < ε
  change ∀ n ≥ N3, |y n - L| < ε / 2 at hN3
  rewrite [show |x n - L| = |(x n - y n) + (y n - L)| by ring_nf]
  have f1 : |(x n - y n) + (y n - L)| ≤
      |(x n - y n)| + |(y n - L)| := by apply abs_add
  specialize hN n (N n) (by bound)
  change |y n - x n| < 1 / (n + 1) at hN
  rewrite [show |y n - x n| = |x n - y n| by apply abs_sub_comm] at hN
  specialize hN3 n (by bound)
  field_simp at hN2
  have hn' : (N1 : R) + N2 + N3 ≤ n := by norm_cast
  have f2' : (N2 : R) ≤ n := by bound
  have f2'' : (N2 : R) * ε ≤ n * ε := by bound
  have f2 : (1 : R) / (n + 1) < ε / 2 := by field_simp; bound
  linarith [f1, f2, hN, hN3]
\end{minted}

\section{Comprendiendo la Demostración}

Desempaquemos los pasos clave en lenguaje natural.

\subsection{La Construcción de $y$}

Usamos la táctica \texttt{choose} para extraer, para cada fila $n$, un índice $N(n)$
tal que la $n$-ésima fila ha convergido hasta dentro de $1/(n+1)$ de su límite en
el índice $N(n)$.

Luego definimos $y(n) = q_{n,N(n)}$: tomamos un elemento de cada fila, elegido en
un punto donde esa fila ya ha convergido suficientemente.

\subsection{Por qué $y$ es de Cauchy}

La demostración de que $y$ es de Cauchy usa una ingeniosa descomposición en tres
términos:
\[
  |y_m - y_n| \leq |y_m - x_m| + |x_m - x_n| + |x_n - y_n|
\]
El término del medio es pequeño porque $x$ es de Cauchy (por hipótesis). El
primero y el tercero son pequeños porque elegimos $y_m$ e $y_n$ en puntos donde
las filas $m$ y $n$ ya habían convergido a sus límites $x_m$ y $x_n$.

Pero hay una sutileza: necesitamos que estos términos sean pequeños
\emph{simultáneamente} para $m$ y $n$ grandes. Aquí es donde entra la propiedad
arquimediana. Necesitamos tanto $1/(m+1) < \varepsilon/3$ como
$1/(n+1) < \varepsilon/3$. Eligiendo $m, n$ suficientemente grandes (concretamente,
$m, n \geq N_2$ donde $N_2$ proviene de la propiedad arquimediana), garantizamos
que ambas desigualdades se cumplen.

\subsection{Por qué $x_n \to L$}

Una vez que sabemos que $y$ es de Cauchy, representa algún número real
$L = \texttt{Real\_of\_CauSeq}\ \mathit{hy}$.
Para demostrar que $x_n \to L$, usamos de nuevo una descomposición en dos términos:
\[
  |x_n - L| \leq |x_n - y_n| + |y_n - L|
\]
El primer término es pequeño porque $y_n$ fue elegido de la fila $n$ en un punto
donde la fila $n$ había convergido a $x_n$. El segundo término es pequeño porque
$y$ converge a $L$ (por \texttt{SeqLim\_of\_Real\_of\_Cau}).

\section{La Importancia de la Completitud}

La completitud de los números reales es lo que hace funcionar el análisis real.
He aquí por qué es importante:

\subsection{Clausura Bajo Límites}

Los números reales son cerrados bajo la operación de tomar límites: si se tiene una
sucesión de Cauchy de números reales, su límite es también un número real. No es
necesario seguir expandiendo hacia sistemas numéricos más grandes.

Esto contrasta marcadamente con los números racionales. La sucesión
$a_n = (1 + 1/n)^n$ es de Cauchy en $\mathbb{Q}$, pero su límite $e$ no es
racional. Para capturar este límite, necesitamos construir $\mathbb{R}$.

\subsection{Teoremas de Punto Fijo}

Muchos teoremas de punto fijo dependen de la completitud. Por ejemplo, el Teorema
del Punto Fijo de Banach establece que una contracción sobre un espacio métrico
completo tiene un único punto fijo, que se encuentra iterando desde cualquier punto
inicial.

\subsection{Convergencia de Series}

La completitud nos permite trabajar con series infinitas. Una serie
$\sum a_n$ converge si y solo si sus sumas parciales forman una sucesión de Cauchy.
Sin completitud, no podríamos garantizar que las sucesiones de Cauchy converjan.

\subsection{Construcción de Funciones}

Muchas funciones importantes se definen como límites de sucesiones:
\begin{itemize}
  \item La función exponencial: $e^x = \lim_{n\to\infty}(1 + x/n)^n$
  \item Las funciones trigonométricas: definidas mediante series de potencias
  \item Funciones especiales: funciones de Bessel, función gamma, etc.
\end{itemize}
Sin completitud, estos límites podrían no existir en el espacio con el que estamos
trabajando.

\section{El Proceso de Completación}

La construcción que acabamos de completar ilustra un principio general en
matemáticas: la \emph{completación} de un espacio.

\subsection{El Patrón General}

Comenzamos con un espacio incompleto $X$ (como $\mathbb{Q}$) dotado de una noción
de distancia. Para completarlo:
\begin{enumerate}
  \item Consideramos todas las sucesiones de Cauchy en $X$.
  \item Definimos una relación de equivalencia: dos sucesiones de Cauchy son
        equivalentes si su diferencia converge a cero.
  \item El espacio completado $\overline{X}$ es el conjunto de clases de
        equivalencia.
  \item Definimos la distancia en $\overline{X}$ usando representantes de las
        clases de equivalencia.
\end{enumerate}

\subsection{Propiedades Clave de la Completación}

\textbf{Universalidad:} La completación es única (salvo isomorfismo). Todo espacio
completo que contenga a $X$ como subconjunto denso debe ser isomorfo a
$\overline{X}$.

\textbf{Preservación:} Las operaciones algebraicas (adición, multiplicación) y las
propiedades topológicas (continuidad) se extienden de $X$ a $\overline{X}$.

\textbf{Completitud:} ¡El espacio completado $\overline{X}$ es completo! Esto es
exactamente lo que demostramos: las sucesiones de Cauchy en $\mathbb{R}$ convergen
en $\mathbb{R}$.

\subsection{Ejemplos en Matemáticas}

\textbf{De $\mathbb{Q}$ a $\mathbb{R}$:} Los números reales son la completación de
los racionales respecto al valor absoluto usual.

\textbf{Números $p$-ádicos:} Usando una noción diferente de distancia (el valor
absoluto $p$-ádico), podemos completar $\mathbb{Q}$ para obtener los números
$p$-ádicos $\mathbb{Q}_p$, que tienen propiedades muy distintas de $\mathbb{R}$.

\textbf{Espacios de funciones:} El espacio $L^2([0,1])$ de funciones de cuadrado
integrable es la completación del espacio de funciones continuas bajo la norma
$L^2$.

\textbf{Completación de espacios métricos:} Cualquier espacio métrico puede
completarse hasta obtener un espacio métrico completo.

\section{El Argumento Diagonal}

La construcción de $y$ eligiendo elementos $y_n = q_{n,N(n)}$ de diferentes filas
es una forma de argumento diagonal.

\section{Reflexiones Filosóficas}

La completitud de los números reales representa un profundo logro matemático: la
construcción de un sistema numérico donde toda sucesión de Cauchy converge.

\subsection{El Precio de la Completitud}

Para lograr la completitud, tuvimos que abandonar el mundo de los racionales, donde
todo número puede escribirse como una expresión finita $p/q$, y entrar en el mundo
de los reales, donde la mayoría de los números (como $\pi$, $e$, $\sqrt{2}$)
requieren expansiones decimales infinitas.

Este es el compromiso fundamental: la completitud exige aceptar los procesos
infinitos como objetos completados.

\subsection{El Beneficio de la Completitud}

A cambio, obtenemos un sistema numérico donde:
\begin{itemize}
  \item Toda sucesión monótona y acotada converge.
  \item Toda sucesión de Cauchy converge.
  \item Las funciones continuas en intervalos cerrados alcanzan sus extremos.
  \item El Teorema del Valor Intermedio se cumple.
  \item Las series de potencias definen funciones.
\end{itemize}
Estas propiedades hacen posible el análisis. Sin completitud, nos encontraríamos
constantemente con sucesiones que «deberían» converger pero no lo hacen, funciones
que «deberían» tener ceros pero no los tienen, y teoremas que «deberían» ser válidos
pero fallan.

%%% Local Variables:
%%% mode: LaTeX
%%% TeX-master: "Analisis_en_Lean"
%%% End:
